% This file is part of the Open Source project 'proTirone'
% (c) 2025 Karsten Reincke (https://github.com/protirone/protico.ltx)
% It is distributed under the terms of the creative commons license
% CC-BY-4.0 (= https://creativecommons.org/licenses/by/4.0/)

\selectlanguage{ngerman}

{
\usebackgroundtemplate{\includegraphics[width=\paperwidth]{\imgLf/beach-steps-725480-pxh.png}}
\begin{frame}  
  \frametitle{LF11c:01:Prozess:Dokumentation}

  \vspace{6cm}
  \begin{flushright}
    \textcolor{white}{Was sind \\
    \textit{Prozesse}?}
  \end{flushright}
\end{frame}
}

\begin{frame}[fragile]
  \frametitle{LF11c:01:Prozess:Dokumentation:cat|tr|sed}
  \begin{columns}
    \column{0.5\textwidth}
      \begin{block}{Beispielprozess}
        \begin{verbatim}

cat MIT.dirty.txt |\
   sed 's/  */ /g' |\
   tr -d '\r' |\
   tr '\n' '$' |\
   sed 's/\$\$/\n\n/g' |\
   tr -d '$' |\
   tr 'a' 'a' |\
   sed 's/  */ /g' \
 >> x.txt
        \end{verbatim}
      \end{block}
    \column{0.5\textwidth} 
    \begin{footnotesize}
    \begin{itemize}
      \item \texttt{cat} :- liest Daten aus Datei und gibt sie unverändert nach stdout aus.
      \item \texttt{sed} :- (= Stream-Editor) = liest Daten von \texttt{stdin}, ersetzt das zwischen \texttt{/ /} durch das danach \texttt{/ /} und gibt das Ergebnis nach \texttt{stdout} aus.
      \item \texttt{tr} :- liest Daten von \texttt{stdin}, ersetzt alle Zeichen 'A' durch 'B' bzw. löscht 'A' und schreibt nach \texttt{stdout}
      \item \texttt{>} :- liest von einem IO-Channel (hier \texttt{stdout}) und lenkt die in eine Datei um.
    \end{itemize}
  \end{footnotesize}
  \end{columns}

\end{frame}

\begin{frame}
  \frametitle{LF11c:01::cat|sed|tr:Aktivitätsdiagramm}
  \begin{center}
    \includegraphics[height=7.5cm]{\imgLf/cat-sed-tr-solution-act.png}
  \end{center}
\end{frame}

\begin{frame}
  \frametitle{LF11c:01::cat|sed|tr:Flussdiagramm}
  \begin{center}
    \includegraphics[height=7.2cm]{\imgLf/cat-sed-tr-solution-flow.png}
  \end{center}
\end{frame}

\begin{frame}[fragile] 
  \frametitle{LF11c:01::BPMN \& (e)EPK:Einkaufen}
  \begin{center}
    \includegraphics[width=10cm]{\imgLf/ldk-stores.png}
  \end{center}

  \vspace{0.4cm}
  \begin{center}
  \begin{minipage}{0.5\textwidth}
    \textbf{Meine Einkaufsregeln im Kern}

    \begin{minted}[fontsize=\small,linenos=true]{python}
if (reinigung): in_Herborn(Reinigung,DM,Rewe)
elif (drogerie): in_Erda(Rossmann,Edeka)
else: in_Niederweidbach(Rewe) 
    \end{minted}
  \end{minipage}
  \end{center}

\end{frame}

\begin{frame}  
  \frametitle{LF11c:01::Einkaufen:BPMN-Diagramm}
  \begin{center}
    \includegraphics[width=10cm]{\imgGl/coming-soon-118519-pxh.png}
  \end{center}
\end{frame}

\begin{frame}  
  \frametitle{LF11c:01::Einkaufen:EPK-Diagramm}
  \begin{center}
    \includegraphics[width=10cm]{\imgLf/shopping-epk.png}
  \end{center}
\end{frame}

\begin{frame}  
  \frametitle{LF11c:01::BPMN \& (e)EPK:Switchi}
  \begin{center}
    \includegraphics[width=10cm]{\imgLf/switchi.png}
  \end{center}
\end{frame}

\begin{frame}
  \frametitle{LF11c:01::Switchi:BPMN Diagramm}
  \begin{center}
    \includegraphics[height=7.8cm]{\imgLf/switchi-bpmn.png}
  \end{center}
\end{frame}

\begin{frame}
  \frametitle{LF11c:01::Switchi:eEPK Diagramm}
  \begin{center}
    \includegraphics[height=7.8cm]{\imgLf/switchi-epc.png}
  \end{center}
\end{frame}